\chapter*{Abstract}
\addcontentsline{toc}{chapter}{Abstract}

The six artist's ledgers of Edward and Josephine Hopper at the Whitney Museum of American Art, 504 digitised sheets kept from 1913 to 1967, and four of Josephine Hopper's notebooks from the Sanborn Hopper Archive were transcribed in September 2026 by a language model, sheet by sheet against the museum's photographs, into a restricted LaTeX with a critical apparatus that marks what was read, what was doubted, what could not be read and whose hand wrote it. This note reads that transcription whole. It describes the ten books and the division of labour between the two hands; follows each ledger through its prices, buyers, exhibitions and descriptions; reads the notebooks as the other half of the same working life; derives from the leaves a chronology of the painter's materials and a census of his formats; and adds up, under stated rules, what the books say was charged and received over fifty-four years. Every quotation is cited by leaf and by the museum's reference and hyperlinked to the transcribed sheet beside its photograph; every figure is a floor under a published rule. The note argues that the edition's contribution is a condition rather than a finding, that any sentence read from these books can now be checked against the hand it came out of on one screen, and it presents itself as an experiment in that method, made with the Hopper corpus as its occasion and with the model's reading, still unchecked by a person, left in view. The methodology chapter sets out the pipeline, from the harvested inventory to the accounts readers, and the working method: Claude Code with Opus 5 and Fable 5.1, a GitHub repository as workbench, an interface built for checking, and parallel sessions by remote control.

\vspace{1em}\noindent\textsc{Keywords.} Edward Hopper; Josephine Nivison Hopper; artist's ledgers; transcription; critical apparatus; TEI; machine reading; provenance; art market; digital humanities.

\vspace{1em}\noindent\textsc{Data availability.} The transcriptions, the derived data, the site and the sources of this note are in the public repository at \url{https://github.com/Commutator-IO/hopper-project}, the site at \url{https://hopper.commutator.io}; the transcriptions are also exported as TEI P5. A Zenodo deposit with a DOI is planned and is tracked as issue 16 of that repository; until it exists, the repository at a named commit is the citable form. The photographs are the Whitney Museum's and are served from its own ResourceSpace; none is stored by the edition.
